% ---------------------------------------------------------------------------
% config/formatting.tex — Formato académico del documento
% ---------------------------------------------------------------------------

% Márgenes académicos
\geometry{
    left=3cm,
    right=2.5cm,
    top=3cm,
    bottom=3cm,
    letterpaper
}

% Interlineado 1,5; sin sangría de párrafo (sin tab) y con espacio entre
% párrafos para que se distingan visualmente.
\onehalfspacing
\setlength{\parindent}{0pt}
\setlength{\parskip}{6pt}
\setlength{\headheight}{16pt}

% Tolerancia de línea (evita overfull hbox)
\pretolerance=1000
\tolerance=2000
\emergencystretch=4em
\hyphenpenalty=10000
\exhyphenpenalty=10000

% Control de flotantes
\raggedbottom
\renewcommand{\topfraction}{0.85}
\renewcommand{\bottomfraction}{0.85}
\renewcommand{\textfraction}{0.1}
\renewcommand{\floatpagefraction}{0.7}

% Hipervínculos
\hypersetup{
    colorlinks=false,
    pdfborder={0 0 0},
    linkcolor=black,
    urlcolor=black,
    citecolor=black
}

% Columnas tabulares controladas (heredadas del anteproyecto)
\newcolumntype{Y}{>{\RaggedRight\arraybackslash\hspace{0pt}}X}
\newcolumntype{L}[1]{>{\RaggedRight\arraybackslash\hspace{0pt}}p{#1}}
\newcolumntype{R}[1]{>{\RaggedLeft\arraybackslash\hspace{0pt}}p{#1}}

% Nombres de índices en español
\addto\captionsspanish{
    \renewcommand{\contentsname}{Índice General}
    \renewcommand{\listfigurename}{Índice de Figuras}
    \renewcommand{\listtablename}{Índice de Tablas}
}

% Títulos de secciones numerados (formato limpio)
% El capítulo arranca en página nueva; diseño block con líneas arriba/abajo,
% fuente 12/14 en negrita y mayúsculas (estilo Google Docs del snippet).
% Las líneas usan el mismo gris que la regla del pie de página.
% -20pt reduce el espacio antes del título; 20pt tras la línea inferior.
\titleformat{\chapter}[block]
  {{\color{gray}\titlerule}\vspace{1ex}\fontsize{12}{14}\bfseries}
  {\thechapter}
  {1em}
  {\MakeUppercase}
  [\vspace{1ex}{\color{gray}\titlerule}]
\titlespacing*{\chapter}{0pt}{-20pt}{20pt}

% Capítulos sin número (Resumen, Abstract, Índices): centrados, tamaño 16,
% texto normal (sin mayúsculas) y SIN las líneas horizontales.
\titleformat{name=\chapter,numberless}[block]
  {\centering\normalfont\fontsize{16}{20}\bfseries}
  {}
  {0pt}
  {}
  {}
\titlespacing*{name=\chapter,numberless}{0pt}{-20pt}{20pt}

\titleformat{\section}
  {\Needspace{8\baselineskip}\normalfont\large\bfseries}
  {\thesection.}
  {0.6em}
  {}

\titleformat{\subsection}
  {\Needspace{6\baselineskip}\normalfont\normalsize\bfseries}
  {\thesubsection.}
  {0.6em}
  {}

\titleformat{\subsubsection}
  {\Needspace{5\baselineskip}\normalfont\normalsize\bfseries}
  {\thesubsubsection.}
  {0.6em}
  {}

\titleformat{\paragraph}
  {\normalfont\normalsize\bfseries}
  {}
  {0em}
  {}

\titlespacing*{\section}{0pt}{0.45cm}{0.25cm}
\titlespacing*{\subsection}{0pt}{0.45cm}{0.2cm}
\titlespacing*{\subsubsection}{0pt}{0.35cm}{0.15cm}
\titlespacing*{\paragraph}{0pt}{0.3cm}{0.15cm}

% Leyendas de tablas y figuras
\captionsetup{
    font=small,
    labelfont=bf,
    labelsep=colon,
    justification=centering
}

% Listas
\setlist[itemize]{
    leftmargin=1.2cm,
    itemsep=1pt,
    topsep=3pt
}

\setlist[enumerate]{
    leftmargin=1.2cm,
    itemsep=2pt,
    topsep=3pt
}

% Colores UBB
\definecolor{ubbblue}{RGB}{25,75,140}
\definecolor{softblue}{RGB}{225,238,250}
\definecolor{softgreen}{RGB}{228,242,225}
\definecolor{softyellow}{RGB}{255,246,214}
\definecolor{softred}{RGB}{250,226,226}
\definecolor{softgray}{RGB}{245,245,245}
\definecolor{linegray}{RGB}{90,90,90}
\definecolor{deepblue}{RGB}{20,58,112}
\definecolor{tealsoft}{RGB}{218,242,239}
\definecolor{lavendersoft}{RGB}{236,232,250}
\definecolor{screenbg}{RGB}{248,250,252}
\definecolor{bordergray}{RGB}{120,130,145}

% Encabezado y pie de página (diseño limpio: sin encabezado, número en gris
% con regla de pie)
\pagestyle{fancy}
\fancyhf{}
\fancyhead{}
\fancyfoot{}
% Barrita: raise -0.13cm deja su punta casi tocando la línea horizontal del pie
% y 0.62cm de alto la alinea con la altura de la cifra.
\fancyfoot[L]{\color{gray}\fontsize{12}{14}\selectfont\textbf{\thepage} \hspace{0.2cm}\smash{\rule[-0.13cm]{1.5pt}{0.62cm}}}
\renewcommand{\headrulewidth}{0pt}
\renewcommand{\footrulewidth}{1pt}
\renewcommand{\footrule}{\hbox to\headwidth{\color{gray}\leaders\hrule height \footrulewidth\hfill}}

% Las páginas de inicio de capítulo usan \thispagestyle{plain}; le damos el
% mismo pie de página para que el número gris aparezca también ahí.
\fancypagestyle{plain}{
  \fancyhf{}
  \fancyfoot[L]{\color{gray}\fontsize{12}{14}\selectfont\textbf{\thepage} \hspace{0.2cm}\smash{\rule[-0.13cm]{1.5pt}{0.62cm}}}
  \renewcommand{\headrulewidth}{0pt}
  \renewcommand{\footrulewidth}{1pt}
  \renewcommand{\footrule}{\hbox to\headwidth{\color{gray}\leaders\hrule height \footrulewidth\hfill}}
}

% Evitar viudas y huérfanas
\widowpenalty=10000
\clubpenalty=10000
\interlinepenalty=200

% Estilo de código: texto en la fuente normal del documento, sin monoespaciada.
% columns=fullflexible es obligatorio porque, sin él, listings asume columnas
% de ancho fijo y la sangría del contenido se destruiría con una fuente proporcional.
\lstset{
    basicstyle=\normalfont\footnotesize,
    columns=fullflexible,
    breaklines=true,
    frame=single,
    framerule=0.3pt,
    rulecolor=\color{linegray},
    backgroundcolor=\color{softgray},
    aboveskip=8pt,
    belowskip=8pt,
    showstringspaces=false
}
% ---------------------------------------------------------------------------
% Estilo unificado de tablas
% ---------------------------------------------------------------------------
\definecolor{tablahead}{RGB}{217,217,217}  % encabezado (gris)
\definecolor{tablazebra}{RGB}{255,255,255} % filas alternadas (desactivadas)
\definecolor{tablatotal}{RGB}{237,237,237} % filas de totales y de grupo
\definecolor{tablaborde}{RGB}{90,90,90}    % bordes de tablas con cuadrícula
\definecolor{tablaetiqueta}{RGB}{242,242,242} % columna de etiquetas (fichas)

\renewcommand{\arraystretch}{1.35}
\setlength{\aboverulesep}{0pt}
\setlength{\belowrulesep}{0pt}
\setlength{\extrarowheight}{1pt}
\setlength{\heavyrulewidth}{1.1pt}
\setlength{\lightrulewidth}{0.5pt}
\setlength{\LTpre}{6pt}
\setlength{\LTpost}{6pt}
\arrayrulecolor{black}

% Columna X proporcional: P{k} ocupa k veces el ancho medio de las columnas X.
\newcolumntype{P}[1]{>{\RaggedRight\arraybackslash\hspace{0pt}\hsize=#1\hsize\linewidth=\hsize}X}
% Celda de encabezado y activación de filas alternadas
\newcommand{\encab}[1]{{\bfseries #1}}
\newcommand{\tablazebra}{\setstretch{1.05}\rowcolors{2}{white}{tablazebra}}
\newcommand{\filaencabezado}{\rowcolor{tablahead}}
\newcommand{\filatotal}{\rowcolor{tablatotal}}
% Fila separadora de grupo dentro de una tabla (por ejemplo, módulo del sistema)
\newcommand{\filagrupo}[2]{\rowcolor{tablatotal}\multicolumn{#1}{l}{\bfseries\itshape #2} \\}
